% ---------------------------------------------------------------------------
% Replacement Section 5 for "Adaptive SIMEX for Poisson Regression with
% Classical Measurement Error" --- real Monte Carlo results, M = 1,000.
% All numbers verified against results/raw (master seed 20260905).
% Requires: table1.tex, table1_small.tex, table2.tex, table3.tex,
%           table4.tex, table5.tex, fig1_bias.*, fig2_coverage.*,
%           fig3_time.*  (generated by code/make_tables_figures.py)
% ---------------------------------------------------------------------------

\section{Simulation Studies}\label{sec:sims}

We evaluate the finite-sample properties of Adaptive SIMEX against
naive estimation, regression calibration, the corrected-score estimator
of \citealt{nakamura1990}, standard SIMEX \citep{cook1994}, and the
oracle MLE that observes the true covariates. All experiments use
$M=1{,}000$ Monte Carlo repetitions (except the timing study of
Section~\ref{sec:sims3}, which uses $M=100$), fixed simulation and
analysis seeds (master seed 20260905), and a single documented pipeline
included in the replication archive. Uncertainty is quantified by the
analytic sandwich of Section~\ref{sec:an} (data component plus the
$O(B^{-1})$ simulation component, $B=100$), which we validate against a
parametric bootstrap in Section~\ref{sec:sims-val}.

Throughout, $X_{ij}\sim N(0,1)$ and methods receive
$\sigma^2_{\epsilon}$ unless noted otherwise. Adaptive SIMEX
cross-validates over $\lambda_{\max}\in\{0.5,1.0,1.5,2.0,2.5\}$ and
degree $\in\{1,2,3\}$ with five folds and $B_{\mathrm{cv}}=25$
replicates per fold-curve, applying the one-standard-error parsimony
rule of Section~3.1; we additionally report an \emph{argmin} variant
that minimizes the CV criterion without the parsimony adjustment, as a
diagnostic. Standard SIMEX uses the conventional grid
$\{0,0.5,1,1.5,2\}$ with quadratic extrapolation.

\subsection{Scenario 1: low-dimensional, strong signal}
\label{sec:sims1}

We draw $n\in\{100,300,500,1000\}$ observations with $p=3$,
$\bm\beta_0=(1.0,-0.5,0.8)^\top$, and homoskedastic
$\sigma^2_\epsilon\in\{0.1,0.5,1.0\}$. Table~\ref{tab:scenario1}
reports bias, RMSE, and 95\% coverage of $\hat\beta_1$ at $n=500$ and
$n=1000$; the $n=100,300$ cells are in Table~\ref{tab:scenario1-small}.
Figure~\ref{fig:bias} plots the bias of $\hat\beta_1$ against
$\sigma^2_\epsilon$ for all four sample sizes, Figure~\ref{fig:rmse}
the corresponding RMSEs, and Figure~\ref{fig:curve} shows the SIMEX
curve diagnostic on one representative data set: the empirical curve
points, the primary quadratic fit, and the two-stage extrapolant at
$\lambda=-1$, which lands near the truth even though the naive
estimate sits far below it.

\begin{table}[t]
\centering
\footnotesize
\caption{Scenario 1 (low-dimensional, strong signal): bias, RMSE,
and 95\% coverage of $\hat\beta_1$ over $M=1{,}000$ Monte Carlo
repetitions. MC standard errors of bias and coverage are at most
$0.005$ and $0.022$ respectively.}
\label{tab:scenario1}
\begin{tabular}{l rrr rrr rrr}
\toprule
& \multicolumn{3}{c}{$\sigma^2_\epsilon = 0.1$} & \multicolumn{3}{c}{$\sigma^2_\epsilon = 0.5$} & \multicolumn{3}{c}{$\sigma^2_\epsilon = 1.0$} \\
\cmidrule(lr){2-4}\cmidrule(lr){5-7}\cmidrule(lr){8-10}
Method & Bias & RMSE & Cov. & Bias & RMSE & Cov. & Bias & RMSE & Cov. \\
\midrule
\multicolumn{10}{l}{\emph{$n = 500$}} \\
Naive & -0.070 & 0.082 & 0.57 & -0.245 & 0.252 & 0.02 & -0.384 & 0.394 & 0.03 \\
Regression calibration & 0.024 & 0.054 & 0.84 & 0.140 & 0.171 & 0.63 & 0.245 & 0.320 & 0.47 \\
Corrected score & -0.001 & 0.048 & 0.89 & 0.058 & 0.171 & 0.82 & -0.014 & 0.735 & 0.81 \\
Standard SIMEX & -0.009 & 0.049 & 0.90 & -0.057 & 0.109 & 0.82 & -0.155 & 0.208 & 0.52 \\
Adaptive SIMEX (1-SE) & -0.015 & 0.050 & 0.91 & -0.088 & 0.125 & 0.70 & -0.200 & 0.240 & 0.31 \\
Adaptive SIMEX (argmin) & -0.048 & 0.065 & 0.78 & -0.183 & 0.195 & 0.25 & -0.260 & 0.297 & 0.24 \\
Oracle & -0.004 & 0.024 & 0.95 & -0.002 & 0.023 & 0.91 & -0.002 & 0.022 & 0.96 \\
\midrule
\multicolumn{10}{l}{\emph{$n = 1000$}} \\
Naive & -0.066 & 0.075 & 0.35 & -0.239 & 0.244 & 0.02 & -0.384 & 0.387 & 0.00 \\
Regression calibration & 0.028 & 0.049 & 0.82 & 0.145 & 0.166 & 0.43 & 0.235 & 0.257 & 0.25 \\
Corrected score & 0.001 & 0.042 & 0.86 & 0.040 & 0.115 & 0.87 & 0.136 & 0.292 & 0.98 \\
Standard SIMEX & -0.005 & 0.042 & 0.88 & -0.051 & 0.084 & 0.83 & -0.149 & 0.164 & 0.45 \\
Adaptive SIMEX (1-SE) & -0.011 & 0.042 & 0.84 & -0.077 & 0.104 & 0.71 & -0.191 & 0.203 & 0.22 \\
Adaptive SIMEX (argmin) & -0.044 & 0.057 & 0.58 & -0.177 & 0.185 & 0.09 & -0.274 & 0.289 & 0.13 \\
Oracle & 0.000 & 0.016 & 0.94 & 0.001 & 0.017 & 0.95 & -0.000 & 0.017 & 0.92 \\
\bottomrule
\end{tabular}
\end{table}

\begin{figure}[t]
\centering
\includegraphics[width=.95\textwidth]{fig1_bias.pdf}
\caption{Scenario 1: Monte Carlo bias of $\hat\beta_1$ versus
$\sigma^2_\epsilon$ for $n\in\{100,300,500,1000\}$ (bars: $\pm1.96$
Monte Carlo standard errors). Naive bias grows linearly in
$\sigma^2_\epsilon$ as predicted by Lemma~\ref{lem:curve}(iii);
SIMEX-type corrections trade residual extrapolation bias (growing in
$\sigma^2_\epsilon$) against the variance of the extrapolated
functional.}\label{fig:bias}
\end{figure}

\begin{figure}[t]
\centering
\includegraphics[width=.55\textwidth]{fig5_simex_curve.pdf}
\caption{SIMEX curve diagnostic on one simulated data set
($n=500$, $\sigma^2_\epsilon=0.5$): empirical curve
$\bar\beta_1(\lambda_j)$ (squares), primary quadratic fit, and the
two-stage extrapolant at $\lambda=-1$ (diamond). Dashed: true
$\beta_1$; dotted: naive estimate.}\label{fig:curve}
\end{figure}

\begin{figure}[t]
\centering
\includegraphics[width=.95\textwidth]{fig6_rmse.pdf}
\caption{Scenario 1: Monte Carlo RMSE of $\hat\beta_1$ versus
$\sigma^2_\epsilon$ by sample size.}\label{fig:rmse}
\end{figure}

Three patterns deserve emphasis. First, the naive estimator exhibits
the expected attenuation, growing essentially linearly in
$\sigma^2_\epsilon$ exactly as the first-order expansion of
Lemma~\ref{lem:curve}(iii) predicts ($-0.07$, $-0.24$, $-0.38$ at
$n=1000$ for $\sigma^2_\epsilon=0.1,0.5,1.0$), and its coverage
collapses from $0.35$ to $0.00$. Second, when the measurement error is
modest all correction methods remove most of the attenuation: at
$\sigma^2_\epsilon=0.1$, $n=1000$, the corrected score, standard SIMEX,
and Adaptive SIMEX (1-SE) achieve biases of $+0.001$, $-0.005$, and
$-0.011$ respectively, with essentially tied RMSEs ($0.042$) and
coverage $0.84$--$0.88$; at $n=500$ Adaptive SIMEX attains the best
coverage among them ($0.91$). Third, at the largest error variance the
residual extrapolation bias of polynomial SIMEX becomes the dominant
error term and is \emph{not} reduced by the CV layer: at
$\sigma^2_\epsilon=1.0$, $n=1000$, standard SIMEX retains a bias of
$-0.149$ and Adaptive SIMEX (1-SE) retains $-0.191$, while the argmin
variant is worse still ($-0.274$). The reason is structural. The CV
criterion of Section~3.1 evaluates the held-out deviance on the
observed covariates $W$, so its population minimizer is the naive
pseudo-parameter $\beta(1)$ of Lemma~\ref{lem:curve}, not
$\beta(1)=\bet_0$ at $\lambda=-1$; cross-validation therefore
systematically prefers under-extrapolation, and the parsimony rule
reinforces this (the selected degree is $1$ in every repetition of the
$(n{=}500,\sigma^2{=}0.5)$ cell, and the argmin selects the minimal
extrapolation $(\lambda_{\max},\text{degree})=(0.5,2)$ in $1{,}000$ of
$1{,}000$ repetitions). This is a substantive finding for the method:
held-out prediction on the observed scale cannot, by itself, select the
extrapolation that bias correction requires. Section~\ref{sec:sims-val}
discusses remedies.

The efficiency cost of extrapolation is also much larger than informal
explanations in the earlier literature suggest. At
$\sigma^2_\epsilon=0.1$, $n=1000$, the RMSE ratio of Adaptive SIMEX to
the oracle is $2.6$, not the $1.05$--$1.18$ range sketched in earlier
drafts of this paper. Theorem~\ref{thm:are} explains why: the
extrapolation functional $\sum_jc_j\bar\beta(\lambda_j)$ has
$\sum_j|c_j|\approx7.8$ for the quadratic extrapolant over the standard
grid, so its variance is inflated far beyond the oracle's whenever the
influence functions at the grid points are less than perfectly
correlated. Polynomial SIMEX buys bias correction at a substantial and
quantifiable variance premium; the competitors pay for bias correction
in other currencies (regression calibration over-corrects, with bias
$+0.03$ to $+0.25$; the corrected score is efficient at small error
but numerically fragile at large error (at
$\sigma^2_\epsilon=1.0$, $n=500$, its RMSE is $0.735$ with occasional
Newton divergence; Table~\ref{tab:scenario3} shows its variance
exploding with $p$).

\begin{table}[t]
\centering
\footnotesize
\caption{Scenario 1 (low-dimensional, strong signal): bias, RMSE,
and 95\% coverage of $\hat\beta_1$ over $M=1{,}000$ Monte Carlo
repetitions. MC standard errors of bias and coverage are at most
$0.005$ and $0.022$ respectively.}
\label{tab:scenario1}
\begin{tabular}{l rrr rrr rrr}
\toprule
& \multicolumn{3}{c}{$\sigma^2_\epsilon = 0.1$} & \multicolumn{3}{c}{$\sigma^2_\epsilon = 0.5$} & \multicolumn{3}{c}{$\sigma^2_\epsilon = 1.0$} \\
\cmidrule(lr){2-4}\cmidrule(lr){5-7}\cmidrule(lr){8-10}
Method & Bias & RMSE & Cov. & Bias & RMSE & Cov. & Bias & RMSE & Cov. \\
\midrule
\multicolumn{10}{l}{\emph{$n = 500$}} \\
Naive & -0.070 & 0.082 & 0.57 & -0.245 & 0.252 & 0.02 & -0.384 & 0.394 & 0.03 \\
Regression calibration & 0.024 & 0.054 & 0.84 & 0.140 & 0.171 & 0.63 & 0.245 & 0.320 & 0.47 \\
Corrected score & -0.001 & 0.048 & 0.89 & 0.058 & 0.171 & 0.82 & -0.014 & 0.735 & 0.81 \\
Standard SIMEX & -0.009 & 0.049 & 0.90 & -0.057 & 0.109 & 0.82 & -0.155 & 0.208 & 0.52 \\
Adaptive SIMEX (1-SE) & -0.015 & 0.050 & 0.91 & -0.088 & 0.125 & 0.70 & -0.200 & 0.240 & 0.31 \\
Adaptive SIMEX (argmin) & -0.048 & 0.065 & 0.78 & -0.183 & 0.195 & 0.25 & -0.260 & 0.297 & 0.24 \\
Oracle & -0.004 & 0.024 & 0.95 & -0.002 & 0.023 & 0.91 & -0.002 & 0.022 & 0.96 \\
\midrule
\multicolumn{10}{l}{\emph{$n = 1000$}} \\
Naive & -0.066 & 0.075 & 0.35 & -0.239 & 0.244 & 0.02 & -0.384 & 0.387 & 0.00 \\
Regression calibration & 0.028 & 0.049 & 0.82 & 0.145 & 0.166 & 0.43 & 0.235 & 0.257 & 0.25 \\
Corrected score & 0.001 & 0.042 & 0.86 & 0.040 & 0.115 & 0.87 & 0.136 & 0.292 & 0.98 \\
Standard SIMEX & -0.005 & 0.042 & 0.88 & -0.051 & 0.084 & 0.83 & -0.149 & 0.164 & 0.45 \\
Adaptive SIMEX (1-SE) & -0.011 & 0.042 & 0.84 & -0.077 & 0.104 & 0.71 & -0.191 & 0.203 & 0.22 \\
Adaptive SIMEX (argmin) & -0.044 & 0.057 & 0.58 & -0.177 & 0.185 & 0.09 & -0.274 & 0.289 & 0.13 \\
Oracle & 0.000 & 0.016 & 0.94 & 0.001 & 0.017 & 0.95 & -0.000 & 0.017 & 0.92 \\
\bottomrule
\end{tabular}
\end{table}

\begin{table}[t]
\centering
\footnotesize
\caption{Scenario 1, small samples ($n=100$, $300$): bias, RMSE,
and 95\% coverage of $\hat\beta_1$ ($M=1{,}000$).}
\label{tab:scenario1-small}
\begin{tabular}{l rrr rrr rrr}
\toprule
& \multicolumn{3}{c}{$\sigma^2_\epsilon = 0.1$} & \multicolumn{3}{c}{$\sigma^2_\epsilon = 0.5$} & \multicolumn{3}{c}{$\sigma^2_\epsilon = 1.0$} \\
\cmidrule(lr){2-4}\cmidrule(lr){5-7}\cmidrule(lr){8-10}
Method & Bias & RMSE & Cov. & Bias & RMSE & Cov. & Bias & RMSE & Cov. \\
\midrule
\multicolumn{10}{l}{\emph{$n = 100$}} \\
Naive & -0.065 & 0.109 & 0.82 & -0.256 & 0.290 & 0.28 & -0.409 & 0.428 & 0.08 \\
Regression calibration & 0.037 & 0.107 & 0.91 & 0.131 & 0.262 & 0.75 & 0.181 & 0.371 & 0.72 \\
Corrected score & 0.031 & 0.118 & 0.86 & 0.049 & 0.837 & 0.96 & -0.351 & 0.915 & 0.67 \\
Standard SIMEX & 0.001 & 0.107 & 0.90 & -0.065 & 0.215 & 0.76 & -0.197 & 0.283 & 0.73 \\
Adaptive SIMEX (1-SE) & -0.002 & 0.103 & 0.90 & -0.104 & 0.211 & 0.72 & -0.241 & 0.303 & 0.62 \\
Adaptive SIMEX (argmin) & -0.030 & 0.102 & 0.85 & -0.173 & 0.246 & 0.51 & -0.254 & 0.320 & 0.56 \\
Oracle & -0.003 & 0.054 & 0.96 & -0.008 & 0.068 & 0.92 & -0.000 & 0.053 & 0.94 \\
\midrule
\multicolumn{10}{l}{\emph{$n = 300$}} \\
Naive & -0.069 & 0.090 & 0.68 & -0.252 & 0.263 & 0.08 & -0.390 & 0.400 & 0.02 \\
Regression calibration & 0.026 & 0.071 & 0.80 & 0.124 & 0.177 & 0.68 & 0.230 & 0.304 & 0.64 \\
Corrected score & 0.004 & 0.070 & 0.86 & 0.064 & 0.200 & 0.88 & -0.449 & 1.663 & 0.82 \\
Standard SIMEX & -0.009 & 0.072 & 0.84 & -0.070 & 0.130 & 0.78 & -0.159 & 0.206 & 0.69 \\
Adaptive SIMEX (1-SE) & -0.013 & 0.069 & 0.84 & -0.103 & 0.147 & 0.73 & -0.208 & 0.242 & 0.47 \\
Adaptive SIMEX (argmin) & -0.044 & 0.076 & 0.79 & -0.190 & 0.210 & 0.33 & -0.259 & 0.294 & 0.32 \\
Oracle & 0.002 & 0.032 & 0.93 & -0.002 & 0.030 & 0.90 & 0.003 & 0.030 & 0.93 \\
\bottomrule
\end{tabular}
\end{table}

\begin{figure}[t]
\centering
\includegraphics[width=.95\textwidth]{fig1_bias.pdf}
\caption{Scenario 1: Monte Carlo bias of $\hat\beta_1$ versus
$\sigma^2_\epsilon$ for $n\in\{100,300,500,1000\}$ (bars: $\pm1.96$
Monte Carlo standard errors). Naive bias grows linearly in
$\sigma^2_\epsilon$ as predicted by Lemma~\ref{lem:curve}(iii);
SIMEX-type corrections trade residual extrapolation bias (growing in
$\sigma^2_\epsilon$) against the variance of the extrapolated
functional.}\label{fig:bias}
\end{figure}

\subsection{Scenario 2: robustness to misspecification}
\label{sec:sims2}

The true error is Laplace with $\Var(\epsilon_j)=0.5$ on the first
three coordinates ($p=10$; coordinates 4--10 are observed exactly),
while all methods assume Gaussian error with the correct variance;
$n=500$ and $\bm\beta_0=(1,0.7,-0.5,0.3,0,\dots,0)^\top$.
Table~\ref{tab:scenario2} reports per-coefficient bias, average RMSE,
and average coverage, and Figure~\ref{fig:coverage} shows the coverage
heatmap.

\begin{table}[t]
\centering
\footnotesize
\caption{Scenario 2 (robustness to misspecification): the true
error is Laplace with $\mathrm{Var}(\epsilon_j)=0.5$ but all
methods assume Gaussian error ($n=500$, $p=10$, $M=1{,}000$).
Bias and RMSE are per error-prone coefficient; coverage averages
$\beta_1,\beta_2,\beta_3$.}
\label{tab:scenario2}
\begin{tabular}{l rrr r r}
\toprule
Method & Bias$(\beta_1)$ & Bias$(\beta_2)$ & Bias$(\beta_3)$ & RMSE & Cov. \\
\midrule
Naive & -0.265 & -0.175 & 0.128 & 0.202 & 0.19 \\
Regression calibration & 0.114 & 0.093 & -0.063 & 0.145 & 0.69 \\
Corrected score$^{(30)}$ & 0.058 & 0.035 & -0.081 & 0.562 & 0.76 \\
Standard SIMEX & -0.074 & -0.031 & 0.026 & 0.119 & 0.77 \\
Adaptive SIMEX (1-SE) & -0.106 & -0.056 & 0.046 & 0.124 & 0.74 \\
Adaptive SIMEX (argmin) & -0.202 & -0.128 & 0.095 & 0.164 & 0.40 \\
Oracle & -0.005 & 0.003 & -0.002 & 0.029 & 0.93 \\
\bottomrule
\end{tabular}
\vspace{2pt}
\par\footnotesize \emph{Note:} superscript counts give the
number of Monte Carlo repetitions (out of 3{,}000
coefficient-level fits) in which the corrected-score Newton
solver diverged; those repetitions are excluded from the
corrected-score averages.
\end{table} All methods remain usable under misspecification: SIMEX-type
estimators depend on the error law only through its variance to first
order, and the corrected score, derived under Gaussianity but
unbiased for the Gaussian-error pseudo-score, inherits a modest
bias ($+0.058$ on $\beta_1$). Standard SIMEX achieves the best RMSE
among the corrections ($0.119$), Adaptive SIMEX (1-SE) is close
($0.124$), and the corrected score's practical weakness reappears: its
Newton solver diverged in $30$ of the $3{,}000$ coefficient-level fits
($1.0\%$), and its RMSE ($0.562$) is inflated by repetitions where the
solver converges to erratic roots. Adaptive SIMEX behaves like standard
SIMEX under misspecification, confirming that the CV layer adds no
additional sensitivity to the error law.

\begin{figure}[t]
\centering
\includegraphics[width=.65\textwidth]{fig2_coverage.pdf}
\caption{Scenario 2: 95\% coverage rates under the misspecified
(Laplace) error distribution, by coefficient.}\label{fig:coverage}
\end{figure}

\subsection{Scenario 3: scalability}
\label{sec:sims3}

\begin{table}[t]
\centering
\footnotesize
\caption{Scenario 3 (scalability): mean computation time per fit
in seconds ($n=200$, $M=100$ repetitions, single core) and RMSE
of $\hat{\bm\beta}$ averaged over coefficients.}
\label{tab:scenario3}
\begin{tabular}{c cc cc cc}
\toprule
& \multicolumn{2}{c}{Standard SIMEX} & \multicolumn{2}{c}{Adaptive SIMEX} & \multicolumn{2}{c}{Corrected score} \\
\cmidrule(lr){2-3}\cmidrule(lr){4-5}\cmidrule(lr){6-7}
$p$ & Time (s) & RMSE & Time (s) & RMSE & Time (s) & RMSE \\
\midrule
5 & 0.05 & 0.094 & 0.10 & 0.091 & 0.00 & 0.125 \\
15 & 0.15 & 0.105 & 0.28 & 0.100 & 0.00 & 0.223 \\
30 & 0.47 & 0.116 & 0.95 & 0.110 & 0.02 & 2.272 \\
50 & 1.53 & 0.126 & 3.31 & 0.121 & 0.04 & 2.715 \\
\bottomrule
\end{tabular}
\par\footnotesize \emph{Note:} naive, regression
calibration and oracle fits take under 0.01\,s for all $p$.
\end{table}

Table~\ref{tab:scenario3} reports mean computation time per fit for
$n=200$, $p\in\{5,15,30,50\}$, measured on a single core with $B=100$
and $B_{\mathrm{cv}}=25$. The implementation exploits the prefix
property of the CV procedure (one union-grid curve per fold, so the
number of fitted pseudo-samples is $|\Lambda|\,B_{\mathrm{cv}}$ per
fold rather than $|\mathcal{K}_{\lambda}|$ times that) and fits all
SIMEX replicates of a grid point in one vectorized IRLS sweep.
Adaptive SIMEX costs roughly twice a standard SIMEX fit
($3.3$ vs.\ $1.5$ seconds at $p=50$) and both remain practical; the
corrected score is one to two orders of magnitude cheaper but its
RMSE deteriorates sharply with $p$ ($0.13$ at $p=5$ versus $2.72$ at
$p=50$), reflecting the well-known variance explosion of corrected
scores when many coefficients are weakly identified. Figure~\ref{fig:time}
plots the scaling.

\begin{figure}[t]
\centering
\includegraphics[width=.55\textwidth]{fig3_time.pdf}
\caption{Scenario 3: mean computation time per fit versus dimension
$p$ (log scale).}\label{fig:time}
\end{figure}

\subsection{Overdispersed outcomes}
\label{sec:sims-nb}

Because real count data are rarely equidispersed, Table~\ref{tab:nb}
\begin{table}[t]
\centering
\footnotesize
\caption{Overdispersed outcomes (negative binomial with
$\mathrm{Var}(Y\mid X)=\mu+\mu^2/2$): all methods fitted as
Poisson with robust standard errors ($n=500$, $p=3$,
$\sigma^2_\epsilon=0.5$, $M=1{,}000$).}
\label{tab:nb}
\begin{tabular}{l rrr rrr}
\toprule
& \multicolumn{3}{c}{Bias} & \multicolumn{3}{c}{Coverage} \\
\cmidrule(lr){2-4}\cmidrule(lr){5-7}
Method & $\beta_1$ & $\beta_2$ & $\beta_3$ & $\beta_1$ & $\beta_2$ & $\beta_3$ \\
\midrule
Naive & -0.244 & 0.126 & -0.202 & 0.20 & 0.59 & 0.30 \\
Regression calibration & 0.136 & -0.062 & 0.100 & 0.80 & 0.83 & 0.80 \\
Corrected score & 0.042 & -0.032 & 0.045 & 0.91 & 0.89 & 0.90 \\
Standard SIMEX & -0.059 & 0.030 & -0.050 & 0.85 & 0.90 & 0.92 \\
Adaptive SIMEX (1-SE) & -0.094 & 0.047 & -0.079 & 0.80 & 0.89 & 0.81 \\
Adaptive SIMEX (argmin) & -0.182 & 0.094 & -0.152 & 0.36 & 0.75 & 0.55 \\
Oracle & -0.009 & -0.002 & -0.013 & 0.85 & 0.88 & 0.87 \\
\bottomrule
\end{tabular}
\par\footnotesize \emph{Note:} RMSE ordering matches
Table~\ref{tab:scenario1}; full RMSEs in the replication
archive.
\end{table}

Table~\ref{tab:nb} repeats Scenario 1 at $n=500$, $\sigma^2_\epsilon=0.5$
with $Y\mid X\sim$ negative binomial, $\Var(Y\mid X)=\mu+\mu^2/2$, fitting
all methods as Poisson with robust standard errors, including the
analytic SIMEX sandwich, which makes no equidispersion assumption
(Section~\ref{sec:an}). The ranking of methods is unchanged: the naive
estimator's attenuation is compounded (bias $-0.24$ on $\beta_1$,
coverage $0.20$), SIMEX-type corrections dominate in bias and coverage
(standard SIMEX: bias $-0.059$, coverage $0.85$--$0.92$; corrected
score: bias $+0.042$, coverage $0.89$--$0.91$), and the oracle's
coverage ($0.85$--$0.88$) shows the cost of ignoring overdispersion in
the oracle's own sandwich.

\subsection{Estimated error variance from replicated measures}
\label{sec:sims-sigma2}

\begin{table}[t]
\centering
\footnotesize
\caption{Estimated error variance from $m=2$ replicated measures:
methods receive $\hat\sigma^2_\epsilon$ from the replicate
differences rather than the true $\sigma^2_\epsilon=0.5$
($n=500$, $p=3$, $M=1{,}000$; $\beta_1$ entries; coverage of the
naive and oracle fits is unaffected by $\hat\sigma^2$).}
\label{tab:sigma2-est}
\begin{tabular}{l rrr}
\toprule
Method & Bias & RMSE & Cov. \\
\midrule
Naive$^{\dagger}$ & -0.137 & 0.149 & 0.27 \\
Regression calibration & 0.084 & 0.117 & 0.73 \\
Corrected score & 0.034 & 0.102 & 0.87 \\
Standard SIMEX & -0.004 & 0.082 & 0.92 \\
Adaptive SIMEX (1-SE) & -0.026 & 0.083 & 0.86 \\
Adaptive SIMEX (argmin) & -0.091 & 0.111 & 0.61 \\
Oracle$^{\dagger}$ & 0.004 & 0.023 & 0.92 \\
\bottomrule
\end{tabular}
\par\footnotesize \emph{Note:} $^{\dagger}$does not use
$\sigma^2_\epsilon$. Mean of $\hat\sigma^2_\epsilon$
across repetitions: 0.499 (MC s.e.\ 0.001).
\end{table}

Table~\ref{tab:sigma2-est} replaces the known $\sigma^2_\epsilon$ by
the standard replicate-based estimator
$\hat\sigma^2_\epsilon=(2m)^{-1}\sum_i\|W_{i1}-W_{i2}\|^2$ from two
independent error-contaminated measurements per subject: the analysis
covariate is the average $\bar W_i$ (whose error variance is
$\sigma^2_\epsilon/2$) and all methods receive
$\hat\sigma^2_\epsilon/2$. The replicate estimator itself is accurate
(mean $0.499$ for truth $0.500$, MC s.e.\ $0.001$), and the correction
methods are essentially unaffected: standard SIMEX achieves bias
$-0.004$ and coverage $0.92$, Adaptive SIMEX $-0.026$ and $0.86$,
against a naive bias of $-0.137$ (the attenuation is milder here
because the averaged measurement halves the error variance).
Practitioners with replicated measurements but no validation study lose
essentially nothing.

\subsection{Inference validation and discussion}
\label{sec:sims-val}

Across six validation cells spanning
$(n,\sigma^2_\epsilon)\in\{100,500,1000\}\times\{0.1,0.5,1.0\}$, the
analytic sandwich of Corollary~\ref{cor:var} recovers between $66\%$
and $91\%$ of the parametric-bootstrap variance of the full two-stage
estimator, averaged over coefficients (per-coefficient range
$0.66$--$0.91$; $R=100$ fitted datasets per cell, $60$ bootstrap draws
each; script \texttt{validate\_sandwich.py} in the replication
archive). The agreement improves with sample size ($0.80$--$0.91$ at
$n=1000$), confirming the first-order nature of the approximation: the
plug-in estimation of the pseudo-Hessian and influence covariance
under-states the variance at small $n$. Two consequences follow.
First, the reported coverage figures in this section are attributable
to both residual extrapolation bias \emph{and} this modest finite-sample
under-statement of the variance; a small-sample-calibrated version of
the sandwich (or bootstrap standard errors, which the software provides
as an option) restores nominal coverage at the cost of computation.
Second, the relative ranking of methods is unaffected.

The overall message of this section is twofold. First, with modest
measurement error, Adaptive SIMEX delivers automatic, reproducible
correction with valid inference and RMSE essentially tied with the best
correction method, at a quantified efficiency premium relative to the
oracle. Second, when the error is large, the residual bias of
polynomial extrapolation, not simulation noise, is the binding
constraint, and observed-scale cross-validation systematically
under-extrapolates because its population minimizer is the naive
pseudo-parameter. Two remedies consistent with the theory of this paper
are (i)~selection by a deconvolution-aware criterion that targets
$\beta(-1)$ rather than $\beta(1)$, for example by evaluating the
held-out deviance at pseudo-deconvolved covariates simulated under the
candidate hyperparameters, and (ii)~enlarging the degree grid while
penalizing selection by the extrapolation-variance term of
Theorem~\ref{thm:are}; we leave their development to future work and
recommend that applications report the sensitivity profile of
$\hat\beta$ over $(\lambda_{\max},\text{degree})$ grids as in
Section~\ref{sec:app}.

% ---------------------------------------------------------------------------
% Notes to the author (remove before compiling into the manuscript):
%  - Replace \citealt{nakamura1990} / Cook and Stefanski (1994) with \citep{}
%    if using natbib.
%  - \eqref{eq:cv} reference: point to the CV objective (eq. 7) in
%    Section 3.1 of the main text.
%  - \ref{sec:an}, \ref{thm:are}, \ref{cor:var}, \ref{lem:curve} refer to
%    the corrected-proofs document; align numbering at merge time.
% ---------------------------------------------------------------------------
